\section*{Chapitre \ref{ch:b1:erosion-transport} --- Érosion et transport}

\begin{solution}{exo:b1:erosion-transport:1}
\qty{0.4}{g.L^{-1}} font \qty{0.4}{kg.m^{-3}}, donc le flux est $500 \times 0.4 =
\qty{200}{kg.s^{-1}}$. Une année de \num{365.25}~jours dure
\qty{3.156e7}{s}, donc le flux annuel est \qty{6.31e9}{kg}, soit
\qty{6.31}{Mt}.
\end{solution}

\begin{solution}{exo:b1:erosion-transport:2}
$w_s = R g d^2/(18\nu) = 1.6548 \times 9.81 \times (\num{30e-6})^2/(18 \times
\num{1.0034e-6}) = \qty{8.09e-4}{m.s^{-1}}$, environ \qty{0.8}{mm} par seconde. Son
nombre de Reynolds est $w_s d/\nu = \num{8.09e-4} \times
\num{30e-6}/\num{1.0034e-6} = 0.024$, bien en dessous de 1~: la loi de Stokes s'applique. Il
coule à travers \qty{2}{m} en $2/\num{8.09e-4} = \qty{2472}{s}$, soit environ
\qty{41}{min}.
\end{solution}

\begin{solution}{exo:b1:erosion-transport:3}
$\tau_b = \rho g h S = 998.2 \times 9.81 \times 0.5 \times \num{4e-4} =
\qty{1.96}{Pa}$. With $(\rho_s - \rho) g = 1651.8 \times 9.81 =
\qty{16204}{N.m^{-3}}$~: pour \qty{1}{mm}, $\theta = 1.96/(16204 \times 0.001) =
0.121$, près de quatre fois le seuil de \num{0.0314}, donc le sable bouge~; pour
\qty{10}{mm}, $\theta = 0.0121$, un cinquième du seuil de \num{0.0556}, donc le
gravier reste. Le ruisseau transporte son sable sur un fond de gravier qui ne bouge pas.
\end{solution}

\begin{solution}{exo:b1:erosion-transport:4}
(a)~$12.6/3.9 = 3.2$~: les rivières transportent environ trois fois plus en suspension
qu'en solution. (b)~Grands réservoirs~: $0.20 \times 14 = \qty{2.8}{Gt}$ par
an~; tous~: $0.26 \times 14 = \qty{3.64}{Gt}$ par an. (c)~Le flux
dissous est de \qty{3.9e15}{g} par an à \qty{120}{g.m^{-3}}, donc l'eau qui le porte
est $\num{3.9e15}/120 = \qty{3.25e13}{m^3}$, environ \qty{32500}{km^3} par an. La
\omterm{def:b1:erosion-transport:load}{charge en suspension} dans cette eau est $\qty{1.26e16}{g}/\qty{3.25e16}{L} =
\qty{0.39}{g.L^{-1}}$.
\end{solution}

\begin{solution}{exo:b1:erosion-transport:5}
(a)~La section a l'aire $10 \times 2 = \qty{20}{m^2}$ et le périmètre
mouillé $10 + 2 \times 2 = \qty{14}{m}$, donc $R_h = 20/14 = \qty{1.43}{m}$.
(b)~$\tau_b = 998.2 \times 9.81 \times 1.43 \times 0.001 = \qty{14.0}{Pa}$,
contre $\qty{19.6}{Pa}$ pour un chenal très large profond de \qty{2}{m}~: le rapport est
$R_h/h = 0.714$, les berges prenant le reste du frottement. (c)~Pour un
chenal rectangulaire de largeur $b$, $R_h = bh/(b + 2h) = h/(1 + 2h/b)$, qui
tend vers $h$ quand $b$ croît.
\end{solution}

\begin{solution}{exo:b1:erosion-transport:6}
(a)~$d_1 = \left(18\nu^2/(Rg)\right)^{1/3} = (18/(Rg))^{1/3}\,\nu^{2/3}$, et $R$
et $g$ ne dépendent pas de la viscosité. (b)~$\nu = \num{1.3059e-3}/999.7 =
\qty{1.306e-6}{m^2.s^{-1}}$ à \qty{10}{\celsius}, contre
\qty{1.0034e-6}{m^2.s^{-1}} à \qty{20}{\celsius}~: le rapport est \num{1.302},
et $1.302^{2/3} = 1.192$, donc $d_1 = 0.104 \times 1.192 = \qty{0.124}{mm}$.
(c)~Dans le domaine de Stokes $w_s$ est inversement proportionnelle à $\nu$~: le grain
chute $1.302$ fois plus lentement, soit environ \qty{30}{\percent}. L'eau froide garde
les sédiments fins en suspension plus longtemps~: une rivière en hiver transporte son silt
un peu plus loin pour le même débit.
\end{solution}

\begin{solution}{exo:b1:erosion-transport:7}
Avec $R g/(18\nu) = \qty{8.99e5}{m^{-1}.s^{-1}}$ et la loi combinée~:
\begin{center}
\begin{tabular}{@{}lccc@{}}
\hline
$d$ & Stokes (\si{m.s^{-1}}) & combinée (\si{m.s^{-1}}) & rapport\\
\hline
\qty{0.1}{mm} & \num{8.99e-3} & \num{7.53e-3} & \num{1.19}\\
\qty{0.5}{mm} & \num{0.225} & \num{0.0711} & \num{3.16}\\
\qty{2}{mm} & \num{3.60} & \num{0.197} & \num{18.3}\\
\hline
\end{tabular}
\end{center}
Le rapport croît avec la taille, parce que le nombre de Reynolds du grain croît~: la loi de Stokes,
valable seulement aux petits nombres de Reynolds, croît comme $d^2$, alors que la traînée du
sillage ralentit les grains grossiers à une vitesse qui ne croît que comme $\sqrt d$.
\end{solution}

\begin{solution}{exo:b1:erosion-transport:8}
(a)~$(Rg/\nu^2)^{1/3} = \qty{25263}{m^{-1}}$. À \qty{0.2}{mm}, $D_* = 5.05$
et $\theta_c = 0.3/7.06 + 0.055(1 - e^{-0.101}) = 0.0425 + 0.0053 =
0.0478$~; $\tau_c = 0.0478 \times 16204 \times \num{2e-4} = \qty{0.155}{Pa}$.
À \qty{20}{mm}, $D_* = 505$ et $\theta_c = 0.3/607 + 0.055(1 - e^{-10.1}) =
0.0005 + 0.0550 = 0.0555$~; $\tau_c = 0.0555 \times 16204 \times 0.02 =
\qty{18.0}{Pa}$. (b)~$h = \tau_c/(\rho g S)$ avec $\rho g S =
\qty{19.58}{Pa.m^{-1}}$~: \qty{7.9}{mm} d'eau déplacent le sable fin, \qty{0.92}{m}
le gravier.
\end{solution}

\begin{solution}{exo:b1:erosion-transport:9}
Le bassin perd \qty{0.8}{Mt} = \qty{8e5}{t} par an, soit $\num{8e5}/2000 =
\qty{400}{t.km^{-2}.\year^{-1}}$. D'après la \cref{prop:b1:erosion-transport:denudation},
$D = \qty{8e8}{kg}/(2650 \times \qty{2e9}{m^2}) = \qty{1.51e-4}{m}$ par an, soit
\qty{151}{m/\Ma}. C'est environ trois fois la médiane des bassins mesurés
au béryllium-10, \qty{54}{m/\Ma}, mais en dessous de leur moyenne de
\qty{218}{m/\Ma}~: un bassin assez rapide, comme en pays raide ou tectoniquement actif.
\end{solution}

\begin{solution}{exo:b1:erosion-transport:10}
(a)~$u_* = \sqrt{g h S} = \sqrt{9.81 \times 3 \times \num{5e-4}} =
\qty{0.1213}{m.s^{-1}}$ et $\tau_b = \rho u_*^2 = \qty{14.7}{Pa}$.
(b)~La suspension demande $w_s < u_* = \qty{0.1213}{m.s^{-1}}$. La loi combinée
donne \qty{0.126}{m.s^{-1}} à \qty{1}{mm}~; résoudre
$w_s(d) = 0.1213$ par essais donne $d = \qty{0.94}{mm}$. (c)~Pour le gravier
$\theta_c \approx 0.0555$, donc $\tau_c = \qty{14.7}{Pa}$ pour $d \approx
14.7/(0.0555 \times 16204) = \qty{16.3}{mm}$~; à cette taille la courbe de seuil donne
bien $\theta_c = 0.0555$. (d)~Les grains plus gros qu'environ \qty{16}{mm} restent sur
le fond et forment un plancher de gravier stable~; le gravier plus fin et le sable plus gros qu'environ
un millimètre voyagent comme \omterm{def:b1:erosion-transport:load}{charge de fond}~; le sable plus fin, le silt et l'argile voyagent en
suspension.
\end{solution}

\begin{solution}{exo:b1:erosion-transport:11}
(a)~$0.30 \times \qty{1000}{m} = \qty{300}{m}$ de roche, en moyenne sur la
région. (b)~Le rebond est $300 \times 2700/3300 = \qty{245}{m}$. (c)~La surface
moyenne n'est abaissée que de $300 - 245 = \qty{55}{m}$. (d)~Les sommets, qui
n'ont pas été érodés, sont portés vers le haut avec toute la région~: ils se tiennent maintenant
\qty{245}{m} plus haut qu'avant. Creuser des vallées peut élever les pics d'une chaîne.
\end{solution}

\begin{solution}{exo:b1:erosion-transport:12}
(a)~$0.9 \times \qty{5e8}{kg}/\qty{1300}{kg.m^{-3}} = \qty{346000}{m^3}$ par an,
soit \qty{0.69}{\percent} du réservoir. (b)~La moitié du réservoir,
\qty{25e6}{m^3}, est remplie en $\num{25e6}/\num{346000} = \num{72}\text{ ans}$ à ce
rythme. (c)~À mesure que le réservoir se remplit, son volume, et donc le temps de séjour de
l'eau, diminue, et des grains plus fins passent au travers~: son efficacité de piégeage baisse.
Pendant ce temps les dépôts plus anciens sont enfouis et compactés, expulsant de l'eau, si bien que
leur masse volumique apparente croît et chaque tonne prend moins de place. Les deux effets ralentissent la
perte de volume à mesure que le réservoir vieillit.
\end{solution}

\begin{solution}{pb:b1:erosion-transport:1}
\textbf{1.} Une année d'eau est $1.2 \times \num{3.156e7} = \qty{3.79e7}{m^3}$~;
étalée sur \qty{8e7}{m^2}, une couche de \qty{0.473}{m}~: un ruissellement de
\qty{473}{mm} par an.

\textbf{2.} $1.2 \times \qty{0.150}{kg.m^{-3}} = \qty{0.18}{kg.s^{-1}}$, soit
\qty{5680}{t} par an.

\textbf{3.} $1.2 \times 0.080 = \qty{0.096}{kg.s^{-1}}$, soit \qty{3030}{t} par
an.

\textbf{4.} La \omterm{def:b1:erosion-transport:load}{charge de fond} est de \qty{568}{t} par an, et le total $5680 + 3030 +
568 = \qty{9278}{t}$ par an.

\textbf{5.} La concentration de sédiments en suspension est élevée pendant les quelques jours
de crue, quand le débit est élevé aussi, et basse le reste du temps. Une
moyenne dans le temps donne à chaque jour le même poids, donc elle est dominée par les nombreux
jours de bas débit et de basse concentration. D'après la
\cref{prop:b1:erosion-transport:flux}, la masse annuelle est
$\int Q\,C\,\mathrm dt$~; la moyenne pondérée par le débit est cette intégrale divisée par
$\int Q\,\mathrm dt$, donc c'est la moyenne pondérée par le débit qui doit être multipliée par
le débit moyen.

\textbf{6.} $3030/9278 = \num{0.33}$, un tiers. Pour les rivières du monde,
$3.9/(12.6 + 3.9) = \num{0.24}$, et moins si la \omterm{def:b1:erosion-transport:load}{charge de fond} était comptée.
La part du bassin est plus élevée~: des collines basses, où une lente érosion mécanique laisse
à la roche le temps de se dissoudre, et du calcaire et de la craie, qui se dissolvent
complètement.

\textbf{7.} $9278/80 = \qty{116}{t.km^{-2}.\year^{-1}}$.

\textbf{8.} $D = \qty{9.278e6}{kg}/(2700 \times \qty{8e7}{m^2}) =
\qty{4.30e-5}{m}$ par an, \qty{43.0}{m/\Ma}.

\textbf{9.} $43.0/54 = 0.80$~: un peu en dessous de la médiane des bassins
mesurés, et cinq fois en dessous de leur moyenne. Un bassin modérément lent, comme on peut s'y attendre
pour de basses collines dans une région calme.

\textbf{10.} Le jaugeage moyenne sur les quelques années de mesure, et dépend
des crues qui s'y sont produites. Le béryllium-10 moyenne sur le temps mis
à éroder les quelques mètres où il est produit~: à \qty{43}{m/\Ma}, chaque mètre
prend \num{23000} ans, donc la moyenne couvre des dizaines de milliers d'années. Ils
peuvent différer à cause de rares grandes crues, absentes d'un enregistrement court, parce que
l'agriculture et les défrichements ont accru l'érosion ces derniers siècles, et parce que
des sédiments sont stockés dans la plaine d'inondation, et repris à elle.

\textbf{11.} Une partie de la \omterm{def:b1:erosion-transport:load}{charge dissoute} ne vient pas de la roche~: une grande part du
bicarbonate porte du carbone pris à l'air ou au sol, et la pluie apporte du sel
marin. Les compter attribuerait au bassin de la roche qu'il n'a pas perdue.

\textbf{12.} $\num{1.5e6}/1.2 = \qty{1.25e6}{s}$, \num{14.5} jours.

\textbf{13.} La loi combinée donne \qty{7.53e-3}{m.s^{-1}}~: le grain tombe de
\qty{8}{m} en \qty{1062}{s}, environ \qty{18}{min}. Il est piégé.

\textbf{14.} $w_s = \num{8.99e5} \times (\num{2e-5})^2 =
\qty{3.60e-4}{m.s^{-1}}$~: \qty{8}{m} en \qty{22200}{s}, environ \num{6.2}~heures.
Piégé.

\textbf{15.} $w_s = \num{8.99e5} \times (\num{1e-6})^2 =
\qty{8.99e-7}{m.s^{-1}}$~: \qty{8}{m} prennent \qty{8.9e6}{s}, \num{103}~jours, sept
fois le temps de séjour. L'essentiel passe à travers le barrage avec l'eau.

\textbf{16.} Le grain doit tomber de \qty{8}{m} en \qty{1.25e6}{s}, à
$\qty{6.4e-6}{m.s^{-1}}$ au moins. D'après la loi de Stokes $d = \sqrt{18\nu
w_s/(Rg)} = \sqrt{\num{6.4e-6}/\num{8.99e5}} = \qty{2.7}{\micro m}$~: une
argile fine. Les grains plus gros sont piégés où qu'ils entrent.

\textbf{17.} $568 + 0.7 \times 5680 = \qty{4544}{t}$ par an, qui remplissent
$\num{4.544e6}/1300 = \qty{3496}{m^3}$ par an, \qty{0.233}{\percent} du
réservoir. La moitié, \qty{7.5e5}{m^3}, est remplie en environ \num{215} ans.

\textbf{18.} $100/43.0 = \qty{2.33}{\Ma}$.

\textbf{19.} $1 - 2700/3300 = 0.18$~: la surface ne descend que de
\qty{18}{\percent} de l'épaisseur enlevée.

\textbf{20.} $100/0.1818 = \qty{550}{m}$ de roche.

\textbf{21.} Toute la région se soulève de $550 \times 2700/3300 = \qty{450}{m}$
de rebond. Les sommets à peine érodés sont portés avec elle, de jusqu'à
\qty{450}{m}, tandis que les vallées s'approfondissent~: le relief croît avant de décliner.

\textbf{22.} Le \omterm{def:b1:erosion-transport:denudation}{taux de dénudation} dépend avant tout des pentes. À mesure que les collines
sont abaissées et leurs pentes s'adoucissent, le taux diminuera, à moins qu'un soulèvement ou un
changement de climat ne les raidisse de nouveau~: le temps calculé au taux d'aujourd'hui est une
borne inférieure.

\textbf{23.} $550/42.95 = 12.80$. Au taux d'aujourd'hui et avec le rebond isostatique,
abaisser les collines de Varden de \qty{100}{m} prend \textbf{12.8}~millions d'années.
\end{solution}
